\documentclass[11pt]{article}
\usepackage[a4paper,margin=1in]{geometry}
\usepackage{booktabs}
\usepackage{graphicx}
\usepackage{hyperref}
\title{Layer Reduction with Fitted-Ternary TerGRASP}
\author{(team members)}
\date{October 2026}
\begin{document}
\maketitle

\section{Report}
\subsection{Mentor feedback}
Insert the feedback and weekly-meeting notes shared with the mentor here.
Work since Presentation 1: BI profiling, hard-removal baseline, documented
failures (affine bridge, naive averaging, weight tying), fitted-ternary
TerGRASP implemented and evaluated, plus SVD low-rank and 2-bit
quantization supporting results.

\subsection{Baselines}
ShortGPT angular-distance pruning (Men et al.), LayerDrop (Fan et al.),
CALM early exits, KD to shallow students, L0 gates. Only ShortGPT was
reproduced as the control because the others require longer training than
the available MPS budget.

\subsection{Datasets}
Salesforce/wikitext wikitext-2-raw-v1 (main), wikitext-103-raw-v1 and
wikitext-2-v1 (comparisons). Same \texttt{join\_nonempty} preprocessing as
Assignment 1; 300k-char train cap, 100k-char validation/test cap, 3k-token
PPL windows. Bias handling: fixed seed, identical windows for all models.

\subsection{Experimental plan}
Models: GPT-2 (124M) and SmolLM-135M (135M). Same SEED=42, same eval
windows, targeted-vs-random removal as the falsification control, rank
sweep $\in\{64,128,256\}$, bitwidth sweep $\in\{8,4,2\}$. Hyperparameters:
bridge rank 32, Adam lr $10^{-2}$ for 80 steps, AbsMean ternary, ridge
$10^{-4}$. Metric: validation perplexity (primary), latency in ms and
parameter count (secondary).

\subsection{Project management}
One-machine MPS setup, no paid GPUs. Work split: dataset and eval code,
TerGRASP implementation, SVD and quantization supporting runs, figures and
report (fill in per-member contribution).

\section{Experiment}
\subsection{Implementation}
\texttt{assignment2.ipynb} contains the full pipeline: dataset preparation,
BI ranking, residual capture, bridge fitting, ternary quantization, block
replacement, and PPL/latency/parameter evaluation.

\subsection{Results}
\begin{table}[h]
\centering
\begin{tabular}{lrrrrr}
\toprule
Model & k=1 & k=2 & k=3 & k=4 & Note \\
\midrule
GPT-2 fitted (this work) & 45.11 & 50.09 & 140.42 & 879.52 & keeps graph depth, ~1.58-bit blocks \\
GPT-2 hard removal & 45.76 & 47.95 & 121.81 & 943.24 & baseline control \\
SmolLM fitted (this work) & 21.71 & 40.49 & 33.05 & 38.49 & \\
SmolLM hard removal & 19.93 & 26.13 & 29.91 & 34.20 & \\
\bottomrule
\end{tabular}
\caption{Validation PPL by number of replaced layers.}
\end{table}

\subsection{Findings}
At $k{=}1$ the fitted ternary bridge matches or slightly beats hard removal
while preserving depth and reducing the replaced blocks to $\sim$1.58
bit/parameter. Higher $k$ tracks hard removal closely. The negative
results (affine bridge repair, naive averaging, weight tying) show that the
layer maps are not affine and not weight-interchangeable; a fitted
low-rank bridge is the minimal fix that preserves PPL.

\section{Conclusion}
A fitted low-rank residual bridge with ternary quantization is a viable
drop-in replacement for the most redundant layers: at $k{=}1$ it is
indistinguishable from removing the layer, but keeps the computation
graph intact and cuts storage. The BI profile doubles as a predictor of
quantization tolerance, giving a single ranking that drives both pruning
and compression decisions.

\end{document}
